\begin{theorem}[泄漏常数与零温响应常数的同一化恒等式]\label{thm:pom-leakage-response-identity}
定义两常数
\[
\lambda_{\mathrm{ZT}}:=\varphi,
\qquad
\chi_{\mathrm{ZT}}:=3-\varphi.
\]
则最大相位族上的互信息极限常数满足严格恒等式
\[
\boxed{\;
I_{\mathrm{even}}=1-\chi_{\mathrm{ZT}}\log_2\lambda_{\mathrm{ZT}},
\qquad
I_{\mathrm{odd}}=\frac{1-\chi_{\mathrm{ZT}}\log_2\lambda_{\mathrm{ZT}}}{2}\; }.
\]
\leanverified{lambdaZT}
\leanverified{chiZT}
\leanverified{iEven}
\leanverified{iOdd}
\leanverified{lambdaZT\_eq}
\leanverified{chiZT\_eq\_three\_sub\_lambdaZT}
\leanverified{chiZT\_explicit}
\leanverified{lambdaZT\_pos}
\leanverified{iOdd\_eq\_iEven\_div\_two}
\leanverified{iEven\_explicit}
\leanverified{iOdd\_explicit}
\leanverified{paper\_pom\_leakage\_response\_identity}
\end{theorem}

\begin{proof}
由定义 \( \chi_{\mathrm{ZT}}\log_2\lambda_{\mathrm{ZT}}=(3-\varphi)\log_2\varphi \)，代入 \(I_{\mathrm{even}}\) 的闭式即得第一式；第二式由 \(I_{\mathrm{odd}}=I_{\mathrm{even}}/2\) 立即推出。证毕。
\end{proof}

\begin{conjecture}[泄漏—响应普适律]\label{conj:pom-leakage-response-universality}
设一类在线归一化核在零温意义下具有 carry-free 骨架 Perron 根 \(\lambda_{\mathrm{ZT}}>1\) 与零温响应常数 \(\chi_{\mathrm{ZT}}>0\)。
则对该类核的“最大相位 \(\Rightarrow\) 隐藏位”机制，互信息在稳定阈值区间内的极限常数仅依赖于 \((\lambda_{\mathrm{ZT}},\chi_{\mathrm{ZT}})\)，并随边界奇偶性分裂为两支：
\[
I_{\mathrm{even}}^{\ast}=1-\chi_{\mathrm{ZT}}\log_2\lambda_{\mathrm{ZT}},
\qquad
I_{\mathrm{odd}}^{\ast}=\frac{1-\chi_{\mathrm{ZT}}\log_2\lambda_{\mathrm{ZT}}}{2}.
\]
其中本节所验证的 \((\lambda_{\mathrm{ZT}},\chi_{\mathrm{ZT}})=(\varphi,3-\varphi)\) 情形可视为零温端点的一致性边界条件。
\end{conjecture}

\endinput

